\section{M7 Proof Obligation Graph}
\label{sec:m7_proof_obligation_graph}
\begingroup
\sloppy
\emergencystretch=3em

This fragment turns the M7 parent target into a fine-grained proof-obligation
graph.  Each node is intended to be small enough for an independent Lean task:
it names the output package or theorem, records the immediate dependencies, and
lists the fields that block the node.  The final assembly should not introduce a
new global constructor; it should fill \texttt{SelectedMixedGlobalInput} from
the existing M6.4 and M6.5 packages.

\begin{definition}[M7 shared index spine]
  \label{obl:m7_shared_index_spine}
  \uses{anc:m6_4_completed_constructor_layer,
    anc:m7_selected_mixed_global_assembly}
  \plannedlean{Stokes.M7SharedIndexSpine}
  This node fixes the common finite active index set, selected partition
  coefficient family, source chart family, and source selected boxes used by all
  later M7 packages.

  \emph{Dependency chain:}
  compact support and partition choices feed the active finset; the active
  finset feeds interior pieces, boundary pieces, artificial-face terms, and
  reconstruction sums.

  \emph{Parallel queue:}
  after this node, the support/smoothness queue, artificial-face queue,
  boundary chart-change queue, and integral-reconstruction queue can start.

  \emph{Blocking fields:}
  exact equality of active finsets across packages; exact equality of the
  coefficient family; common term functions for project-local interior terms,
  artificial-boundary terms, and boundary partition terms.
\end{definition}

\begin{definition}[M7 selected partition witness]
  \label{obl:m7_selected_partition_witness}
  \uses{obl:m7_shared_index_spine,
    anc:m6_4_partition_localization_support_control,
    def:partition_subordinate_selected_boxes}
  \plannedlean{Stokes.m7SelectedPartitionWitness}
  This node instantiates the selected partition of unity for the active compact
  support and exposes the coefficient family consumed by localized forms.

  \emph{Dependency chain:}
  \texttt{SelectedBoxPartitionOfUnity} gives coefficient support, finite
  activity, and sum-one fields; those become the inputs for localized support
  and localized reconstruction.

  \emph{Parallel queue:}
  coefficient support and sum-one reconstruction can be proved independently
  once the same active finset and coefficient function have been fixed.

  \emph{Blocking fields:}
  support containment of each coefficient in the selected box; finite active
  membership lemmas; pointwise sum-one on the compact control set.
\end{definition}

\begin{definition}[M7 compact support to active boxes]
  \label{obl:m7_compact_support_active_boxes}
  \uses{obl:m7_selected_partition_witness,
    anc:m6_4_active_compact_support_to_boxes,
    thm:finite_selected_box_cover_exists}
  \plannedlean{Stokes.m7ActiveCompactBoxes}
  This node builds the active selected and extended box data used by interior
  localized pieces.

  \emph{Dependency chain:}
  compact coordinate images produce selected boxes; selected boxes plus
  localized smoothness neighborhoods produce extended boxes.

  \emph{Parallel queue:}
  target-domain shrinkage, overlap shrinkage, compact coordinate support, and
  smooth-neighborhood construction can be split across workers.

  \emph{Blocking fields:}
  \texttt{tsupport\_subset\_coordSupport}; compactness of coordinate supports;
  lower and upper corners with \(a_i \le b_i\); target and overlap inclusions;
  one open \texttt{ContDiffOn} neighborhood for every active index.
\end{definition}

\begin{lemma}[M7 localized support controls]
  \label{obl:m7_localized_support_controls}
  \uses{obl:m7_selected_partition_witness,
    obl:m7_compact_support_active_boxes,
    anc:m6_4_partition_localization_support_control}
  \plannedlean{Stokes.m7LocalizedSupportControls}
  This node proves that every partition-localized chart representative is
  supported in its selected active box.

  \emph{Dependency chain:}
  coefficient support and original-form support combine to bound the support of
  \(\rho_i\omega\); the bound is converted into
  \texttt{LocalizedSupportControl}.

  \emph{Parallel queue:}
  support of scalar multiplication and support transfer through chart
  representatives can be developed independently.

  \emph{Blocking fields:}
  equality between the chosen coefficient and the selected partition
  coefficient; topological-support or \texttt{tsupport} containment for scalar
  multiplication; chart representative support transfer.
\end{lemma}

\begin{lemma}[M7 localized smoothness witnesses]
  \label{obl:m7_localized_smoothness_witnesses}
  \uses{obl:m7_localized_support_controls,
    anc:m6_4_partition_localization_support_control,
    thm:chartwise_smooth_transport}
  \plannedlean{Stokes.m7LocalizedSmoothnessWitnesses}
  This node supplies the smooth-neighborhood witness needed to upgrade selected
  localized boxes to extended boxes.

  \emph{Dependency chain:}
  scalar smoothness of \(\rho_i\), chartwise smoothness of \(\omega\), and the
  product rule for localized forms feed \texttt{LocalizedSmoothnessData}.

  \emph{Parallel queue:}
  scalar partition smoothness, chartwise transport, and localized product
  smoothness can be proved as separate lemmas.

  \emph{Blocking fields:}
  open source neighborhood; containment of the selected box in that
  neighborhood; \texttt{ContDiffOn} for the localized chart representative.
\end{lemma}

\begin{lemma}[M7 localized exterior derivative on support]
  \label{obl:m7_extderiv_on_support}
  \uses{obl:m7_selected_partition_witness,
    obl:m7_localized_smoothness_witnesses,
    anc:m6_5_global_integral_extderiv_reconstruction,
    lem:extderiv_finite_sum}
  \plannedlean{Stokes.m7ChartwiseExtDerivOnSupport}
  This node proves the support-local equality between the chart representative
  of the exterior derivative of the localized finite sum and the chart
  representative of the exterior derivative of the original form.

  \emph{Dependency chain:}
  finite sum reconstruction gives equality of forms on the control support;
  \texttt{extDeriv} finite-sum and chartwise transport convert it into the
  field required by \texttt{ExtDerivOnSupportData}.

  \emph{Parallel queue:}
  this can run in parallel with artificial-face and boundary chart-change work
  once the selected partition witness is fixed.

  \emph{Blocking fields:}
  exact support set for the global bulk reconstruction; finite-sum
  \texttt{extDeriv} equality; restriction of the equality to the integration
  support.
\end{lemma}

\begin{definition}[M7 localized interior piece family]
  \label{obl:m7_localized_interior_piece_family}
  \uses{obl:m7_localized_support_controls,
    obl:m7_localized_smoothness_witnesses,
    anc:m6_4_interior_pieces_artificial_faces,
    thm:interior_chart_local_stokes}
  \plannedlean{Stokes.m7LocalizedInteriorPieces}
  This node constructs the finite family of localized interior pieces and the
  local project-Stokes equality for each piece.

  \emph{Dependency chain:}
  support control gives selected boxes; smoothness witnesses give extended
  boxes; extended boxes feed \texttt{InteriorLocalStokesData}; the finite family
  feeds the artificial-face queue.

  \emph{Parallel queue:}
  individual active indices can be proved independently, then bundled into
  \texttt{LocalizedInteriorPieces}.

  \emph{Blocking fields:}
  selected box accessors; extended box accessors; local integral term equality;
  alignment between the per-piece local term and the global project-local term.
\end{definition}

\begin{definition}[M7 artificial face labels]
  \label{obl:m7_artificial_face_labels}
  \uses{obl:m7_localized_interior_piece_family,
    anc:m6_4_interior_pieces_artificial_faces}
  \plannedlean{Stokes.m7ArtificialFaceLabels}
  This node defines the flattened artificial-face index type and the signed face
  term attached to each interior localized box.

  \emph{Dependency chain:}
  localized interior boxes determine ordinary cubical faces; boundary faces
  that lie on the real manifold boundary are excluded from this artificial
  queue; the remaining faces become cancellation terms.

  \emph{Parallel queue:}
  face indexing and unsigned face-integral comparison can proceed independently
  of the pairing proof.

  \emph{Blocking fields:}
  face-domain containment; proof that the face is artificial; sign convention
  for lower and upper faces; equality with the boundary term emitted by the
  local Stokes package.
\end{definition}

\begin{lemma}[M7 artificial face pairing]
  \label{obl:m7_artificial_face_pairing}
  \uses{obl:m7_artificial_face_labels,
    anc:m6_4_interior_pieces_artificial_faces}
  \plannedlean{Stokes.m7ArtificialFacePairing}
  This node proves that artificial faces pair by opposite orientation and equal
  unsigned integral.

  \emph{Dependency chain:}
  face labels produce a pairing candidate; geometric equality of paired faces
  gives equality of unsigned terms; cubical orientation gives opposite signs.

  \emph{Parallel queue:}
  fixed-point-freeness, geometric face equality, unsigned term equality, and
  sign opposition can be separate lemmas.

  \emph{Blocking fields:}
  active-finset preservation; involutivity; no fixed points or explicit
  fixed-point-zero proof; paired-domain equality; paired integrand equality;
  opposite signed face coefficients.
\end{lemma}

\begin{theorem}[M7 artificial boundary cancellation]
  \label{obl:m7_artificial_boundary_cancellation}
  \uses{obl:m7_artificial_face_labels,
    obl:m7_artificial_face_pairing,
    anc:m6_4_interior_pieces_artificial_faces}
  \plannedlean{Stokes.m7ArtificialBoundaryCancellation}
  This node produces the cancellation field consumed by the mixed global
  constructor.

  \emph{Dependency chain:}
  the local artificial-boundary term is rewritten as the signed face sum; the
  pairing theorem cancels that sum; the result is the
  \texttt{interiorBoundaryCancellation} field.

  \emph{Parallel queue:}
  this waits only for labels, local-term identification, and pairing; it does
  not need global bulk or boundary reconstruction.

  \emph{Blocking fields:}
  \texttt{artificialCancellation\_active};
  \texttt{artificialCancellation\_pieces};
  \texttt{artificialCancellation\_term}; equality between the local Stokes
  boundary term and the chosen artificial-face sum.
\end{theorem}

\begin{definition}[M7 boundary source piece family]
  \label{obl:m7_boundary_source_piece_family}
  \uses{obl:m7_shared_index_spine,
    anc:m6_4_boundary_pieces_chart_change_fields,
    thm:boundary_chart_local_stokes}
  \plannedlean{Stokes.m7BoundarySourcePieces}
  This node constructs source-side boundary local pieces from selected
  boundary chart boxes.

  \emph{Dependency chain:}
  active boundary indices select source boundary boxes; selected image data
  invokes boundary chart local Stokes; the resulting source terms feed
  chart-change cancellation.

  \emph{Parallel queue:}
  source image data can be proved independently for each boundary active
  index.

  \emph{Blocking fields:}
  boundary-active index set; source lower and upper corners; lower-face image
  data; outward-first orientation convention; equality with the source
  boundary partition term.
\end{definition}

\begin{definition}[M7 boundary target image boxes]
  \label{obl:m7_boundary_target_image_boxes}
  \uses{obl:m7_boundary_source_piece_family,
    anc:m6_4_boundary_pieces_chart_change_fields,
    thm:exists_boundary_local_inverse_data}
  \plannedlean{Stokes.m7BoundaryTargetImageBoxes}
  This node chooses target boundary boxes and local inverse data for the
  source-to-target chart-change comparison.

  \emph{Dependency chain:}
  source boundary boxes and chart overlap data determine target boxes; local
  inverse data identifies the source lower face with the target boundary box.

  \emph{Parallel queue:}
  target box selection, local inverse construction, and target support
  containment can be split by boundary active index.

  \emph{Blocking fields:}
  target chart choices; target lower and upper corners; target selected or
  extended boxes; source-to-target image equality; support containment after
  chart change.
\end{definition}

\begin{lemma}[M7 oriented boundary chart change pieces]
  \label{obl:m7_oriented_boundary_chart_change}
  \uses{obl:m7_boundary_source_piece_family,
    obl:m7_boundary_target_image_boxes,
    anc:m6_4_boundary_pieces_chart_change_fields,
    thm:boundary_chart_integral_invariance}
  \plannedlean{Stokes.m7OrientedBoundaryChartChangePieces}
  This node proves the oriented change-of-variables equality for every active
  boundary piece.

  \emph{Dependency chain:}
  local inverse data supplies the map; boundary chart integral invariance
  supplies the integral equality; orientation comparison supplies the sign.

  \emph{Parallel queue:}
  each boundary active index is independent once source and target boxes are
  fixed.

  \emph{Blocking fields:}
  derivative orientation sign; Jacobian or orientation-map compatibility;
  target-domain containment; pointwise equality of pulled-back boundary
  representatives.
\end{lemma}

\begin{theorem}[M7 boundary chart-change cancellation]
  \label{obl:m7_boundary_chart_change_cancellation}
  \uses{obl:m7_oriented_boundary_chart_change,
    anc:m6_4_boundary_pieces_chart_change_fields}
  \plannedlean{Stokes.m7BoundaryChartChangeCancellation}
  This node produces the finite chart-change cancellation package used by the
  mixed global constructor.

  \emph{Dependency chain:}
  oriented per-piece chart-change equalities are summed over the active
  boundary finset; the source sum is identified with the old boundary term;
  the target sum is identified with the new boundary partition term.

  \emph{Parallel queue:}
  old-term and new-term identifications can be assigned separately after the
  per-piece chart-change lemmas exist.

  \emph{Blocking fields:}
  \texttt{chartChange\_active}; \texttt{chartChange\_pieces};
  \texttt{chartChange\_oldTerm}; \texttt{chartChange\_newTerm};
  \texttt{boundaryPartitionTerm\_eq}.
\end{theorem}

\begin{theorem}[M7 global bulk reconstruction]
  \label{obl:m7_global_bulk_reconstruction}
  \uses{obl:m7_selected_partition_witness,
    obl:m7_localized_interior_piece_family,
    anc:m6_5_global_integral_extderiv_reconstruction,
    thm:global_bulk_eq_local_sum}
  \plannedlean{Stokes.m7GlobalBulkReconstruction}
  This node proves that the represented global bulk integral equals the selected
  sum of local project-bulk terms.

  \emph{Dependency chain:}
  finite additivity over localized forms and partition-sum reconstruction
  reduce the global bulk side to the active local terms.

  \emph{Parallel queue:}
  this can run in parallel with boundary reconstruction after the shared
  coefficient family and local term functions are fixed.

  \emph{Blocking fields:}
  definition or comparison theorem for the represented global bulk integral;
  finite additivity over the active finset; equality between local bulk terms
  and project-local bulk terms; compatibility with the support-local
  \texttt{extDeriv} field.
\end{theorem}

\begin{theorem}[M7 global boundary reconstruction]
  \label{obl:m7_global_boundary_reconstruction}
  \uses{obl:m7_boundary_chart_change_cancellation,
    anc:m6_5_global_integral_extderiv_reconstruction,
    thm:global_boundary_eq_local_sum}
  \plannedlean{Stokes.m7GlobalBoundaryReconstruction}
  This node proves that the represented global boundary integral equals the
  selected boundary partition sum after chart-change normalization.

  \emph{Dependency chain:}
  boundary local pieces give source boundary terms; chart-change cancellation
  rewrites them to target partition terms; finite additivity identifies the
  resulting sum with the chosen global boundary integral.

  \emph{Parallel queue:}
  finite additivity and target-term comparison can proceed separately once the
  target boundary partition term function is fixed.

  \emph{Blocking fields:}
  definition or comparison theorem for the represented global boundary
  integral; finite additivity over boundary active terms; equality between
  target chart-change terms and the global boundary partition term.
\end{theorem}

\begin{definition}[M7 reconstruction data fusion]
  \label{obl:m7_reconstruction_data_fusion}
  \uses{obl:m7_extderiv_on_support,
    obl:m7_global_bulk_reconstruction,
    obl:m7_global_boundary_reconstruction,
    anc:m6_5_global_integral_extderiv_reconstruction}
  \plannedlean{Stokes.m7PartitionReconstructionData}
  This node fuses bulk reconstruction, boundary reconstruction, and
  support-local exterior-derivative reconstruction into the reconstruction
  package expected by \texttt{SelectedMixedGlobalInput}.

  \emph{Dependency chain:}
  bulk reconstruction supplies the left side; boundary reconstruction supplies
  the right side; support-local exterior derivative aligns the form being
  integrated; the wrapper constructors produce \texttt{PartitionReconstructionData}
  or \texttt{BulkIntegralReconstructionData} with boundary fields.

  \emph{Parallel queue:}
  this is a synchronization node, not an analysis node.

  \emph{Blocking fields:}
  identical active finset and coefficient family in all three reconstruction
  inputs; identical local bulk term function; identical boundary partition term
  function; exact equality of represented global integral names.
\end{definition}

\begin{theorem}[M7 selected mixed input fields]
  \label{obl:m7_selected_mixed_input_fields}
  \uses{obl:m7_shared_index_spine,
    obl:m7_reconstruction_data_fusion,
    obl:m7_localized_interior_piece_family,
    obl:m7_artificial_boundary_cancellation,
    obl:m7_boundary_chart_change_cancellation,
    anc:m7_selected_mixed_global_assembly}
  \plannedlean{Stokes.m7SelectedMixedGlobalInput}
  This node fills \texttt{SelectedMixedGlobalInput} from the completed queues.

  \emph{Dependency chain:}
  shared indices align the selected partition, interior pieces, artificial
  cancellation, boundary pieces, chart-change cancellation, and reconstruction
  package; the existing constructor converts the record to
  \texttt{MixedGlobalStokesData}.

  \emph{Parallel queue:}
  this is the final M7 assembly gate.

  \emph{Blocking fields:}
  the eight alignment fields
  \texttt{selectedPartition\_active},
  \texttt{artificialCancellation\_active},
  \texttt{artificialCancellation\_pieces},
  \texttt{artificialCancellation\_term},
  \texttt{chartChange\_active},
  \texttt{chartChange\_pieces},
  \texttt{chartChange\_oldTerm}, and
  \texttt{chartChange\_newTerm}; all should be definitional or short rewrite
  goals if upstream queues chose the same names.
\end{theorem}

\begin{theorem}[M7 compact-support theorem from selected input]
  \label{obl:m7_compact_support_from_selected_input}
  \uses{obl:m7_selected_mixed_input_fields,
    anc:m7_selected_mixed_global_assembly,
    thm:global_compact_support_stokes}
  \plannedlean{Stokes.globalStokes_compactSupport_from_selectedInput}
  This node states the compact-support theorem as an application of
  \texttt{selectedMixedGlobalStokes} to the fully populated selected input.

  \emph{Dependency chain:}
  compact support creates the selected input; the selected input gives the mixed
  global Stokes equality; the theorem exposes the project-facing compact-support
  statement.

  \emph{Parallel queue:}
  no independent work should remain here except theorem-shape adaptation.

  \emph{Blocking fields:}
  theorem statement must use the same global bulk and boundary integral API as
  the reconstruction package; compact-support hypotheses must expose the
  selected partition and compact control set expected by the selected input.
\end{theorem}

\begin{theorem}[M7 compact manifold specialization]
  \label{obl:m7_compact_manifold_specialization}
  \uses{obl:m7_compact_support_from_selected_input,
    thm:compact_manifold_stokes}
  \plannedlean{Stokes.globalStokes_compactManifold}
  This node specializes the compact-support theorem to the standard compact
  oriented smooth manifold-with-boundary statement.

  \emph{Dependency chain:}
  compactness of the manifold gives compact support for the form or a
  compact-support wrapper; the compact-support theorem supplies the final
  equality.

  \emph{Parallel queue:}
  this should wait until the compact-support theorem shape is stable.

  \emph{Blocking fields:}
  compact-support coercion or wrapper; theorem statement matching the chosen
  manifold integral API; any comparison theorem to a future canonical mathlib
  manifold-form integral.
\end{theorem}

\begin{proposition}[M7 critical dependency chain]
  \label{obl:m7_critical_dependency_chain}
  \uses{obl:m7_shared_index_spine,
    obl:m7_selected_partition_witness,
    obl:m7_compact_support_active_boxes,
    obl:m7_localized_support_controls,
    obl:m7_localized_smoothness_witnesses,
    obl:m7_extderiv_on_support,
    obl:m7_localized_interior_piece_family,
    obl:m7_artificial_boundary_cancellation,
    obl:m7_boundary_chart_change_cancellation,
    obl:m7_reconstruction_data_fusion,
    obl:m7_selected_mixed_input_fields,
    obl:m7_compact_support_from_selected_input}
  The critical path is:
  \begin{enumerate}
    \item shared indices to selected partition;
    \item selected partition to active boxes;
    \item active boxes to support and smoothness;
    \item support and smoothness to interior local pieces;
    \item interior local pieces to artificial cancellation;
    \item artificial cancellation to selected mixed input;
    \item selected mixed input to compact-support Stokes.
  \end{enumerate}
  The boundary chart-change queue and the global reconstruction queue attach to
  this path at the shared-index and local-term synchronization points.  They
  should not force redesign of the final constructor layer.
\end{proposition}

\begin{theorem}[M7 parallel worker queue]
  \label{obl:m7_parallel_worker_queue}
  \uses{obl:m7_critical_dependency_chain,
    obl:m7_artificial_face_pairing,
    obl:m7_oriented_boundary_chart_change,
    obl:m7_global_bulk_reconstruction,
    obl:m7_global_boundary_reconstruction}
  The recommended parallel queue is:
  \begin{enumerate}
    \item Index spine and selected partition: fix active labels, coefficients,
      source charts, and term names.
    \item Compact geometry: construct active selected boxes, extended boxes,
      target inclusions, and overlap inclusions.
    \item Localized analysis: prove support control, localized smoothness, and
      support-local exterior derivative equality.
    \item Interior local queue: build localized interior pieces, artificial
      face labels, pairing, and cancellation.
    \item Boundary queue: build source boundary pieces, target image boxes,
      oriented chart-change pieces, and chart-change cancellation.
    \item Reconstruction queue: prove global bulk reconstruction, global
      boundary reconstruction, and fusion into the reconstruction data package.
    \item Final assembly: fill \texttt{SelectedMixedGlobalInput}, call
      \texttt{selectedMixedGlobalStokes}, and expose the compact-support and
      compact-manifold theorem shapes.
  \end{enumerate}
  Queues 2, 3, 4, 5, and 6 can overlap after Queue 1 fixes names.  Queue 7
  should be mostly record-field alignment.
\end{theorem}

\begin{remark}[M7 blocker ledger]
  \label{obl:m7_blocker_ledger}
  \uses{obl:m7_parallel_worker_queue}
  The blockers that should be reported explicitly by future workers are:
  \begin{enumerate}
    \item active finset mismatch;
    \item coefficient family mismatch;
    \item selected or extended box containment not strong enough;
    \item unavailable \texttt{ContDiffOn} neighborhood;
    \item artificial face term not matching the local Stokes boundary term;
    \item missing fixed-point-free pairing or fixed-point-zero proof;
    \item boundary target image data missing;
    \item orientation sign mismatch in boundary chart change;
    \item global bulk or boundary integral API mismatch;
    \item support-local \texttt{extDeriv} equality not stated on the integration
      support.
  \end{enumerate}
  A worker blocked by one of these should return the exact field name and the
  upstream node above that owns it.
\end{remark}

\endgroup
